\documentclass{article}
\usepackage[margin=1in]{geometry}
\usepackage{mathtools, amsfonts, amsthm, xcolor, listings}

\title{Title}
\author{Your name here}

\begin{document}
\maketitle

\section*{Chapter 13}

https://github.com/samlyme/Assignments/tree/cs4080/ch13/challenge

\begin{enumerate}
    \item {
        Lox supports only single inheritance—a class may have a single superclass and that’s the only way to reuse methods across classes. Other languages have explored a variety of ways to more freely reuse and share capabilities across classes: mixins, traits, multiple inheritance, virtual inheritance, extension methods, etc.

If you were to add some feature along these lines to Lox, which would you pick and why? If you’re feeling courageous (and you should be at this point), go ahead and add it.

        I originally set out to implement a trait system, but they don't really 
        ``work'' in dynamically typed languages. Instead, I implemented a way 
        for programmers to extend the funcitonality of a \emph{class} after 
        it has been declared. This is similar to traits, but a bit more loosey
        goosey, without the strict typing system.
        
        All future and \emph{current} instances of the class 
        will receive the new methods. Because everything is bound dynamically 
        anyway, this actually wasn't too difficult to implement.

        Example of the syntax: (the `C' class is defined earlier)
        \begin{verbatim}
var sam = C("sam");
var bob = C("bob");

// error
// sam.added();

// Cursed syntax for introspection
class C + {
    added() {
        print this.name + "this method was added after C was defined";
    }
}

sam.added();
bob.added();
// from: sam
// this method was added after C was defined
// from: bob
// this method was added after C was defined
        \end{verbatim}

        The main difficulties of implementing this is in the Resolver. Because 
        there is no way to statically check if the baseClass of this statement is 
        a subclass or not. I ended up comprimising, letting the super always have 
        access to super, and defering the error to runtime.
        \begin{verbatim}
    @Override
    public Void visitClassPlusStmt(Stmt.ClassPlus stmt) {
        ClassType enclosingClass = currentClass;

        currentClass = ClassType.SUBCLASS; // although not "correct", good enough.
        // No good way to statically determine if this class is a subclass, so we just
        // let the user use SUPER.

        resolve(stmt.baseClass);
        beginScope();
        scopes.peek().put("this", true);

        for (Stmt.Function method : stmt.methods) {
            FunctionType declaration = FunctionType.METHOD;
            if (method.name.lexeme.equals("init")) {
                declaration = FunctionType.INITIALIZER;
            }
            resolveFunction(method, declaration);
        }
        endScope();

        currentClass = enclosingClass;

        return null;
    }
        \end{verbatim}
    }

    \item {
        Take out Lox’s current overriding and super behavior and replace it with BETA’s semantics. In short:

When calling a method on a class, prefer the method highest on the class’s inheritance chain.

Inside the body of a method, a call to inner looks for a method with the same name in the nearest subclass along the inheritance chain between the class containing the inner and the class of this. If there is no matching method, the inner call does nothing.

        I did this by replacing the `findMethod' method inside of the LoxClass
        class with `findMethodRoot' that recursively returns the superclass's 
        method, with the closure modified to have an `inner' variable bound.

        Most other code was just plumbing to make sure the variables end up where 
        they're supposed to.

        \begin{verbatim}
    LoxFunction findMethodRoot(String name) {
        if (superclass == null) {
            if (methods.containsKey(name)) return methods.get(name);
            return null;
        }

        if (methods.containsKey(name)) {
            return superclass.findMethodRoot(name).bindInner(methods.get(name));
        }
        return superclass.findMethodRoot(name);
    }
        \end{verbatim}
    }

    \item {
        In the chapter where I introduced Lox, I challenged you to come up with a couple of features you think the language is missing. Now that you know how to build an interpreter, implement one of those features.

        I made a mistake by wanting to have type checking since it is an 
        extremely difficult feature to implement. Nevertheless, I have given it 
        my best shot, despite the absolute cursed code that was required.

        This feature required such extensive changes, that I have separated from 
        the other challenge branches. You can find the code here:

        https://github.com/samlyme/Assignments/tree/cs4080/ch13/challenge-3

        The type checking is extremely limited. It only tracks the class of variables 
        who were intialized with a constructor of a class. All other types are 
        ignored. Also, the type checking doesn't fully function, since reassigning 
        a variable breaks the invariant held in the Resolver class. There is no 
        way to know what the right hand side of an set expression evaluates to,
        and this type of static analysis will likely require a total refactor 
        of the Resolver class, or an entirely new TypeChecker class. 

        To track what methods a class has access to, we add some 
        code to the visitClassStmt code:
        \begin{verbatim}
        scopes.peek().put("this", true);

        Set<String> methodNames = new HashSet<>();
        classProperties.put(stmt.name.lexeme, methodNames);

        for (Stmt.Function method : stmt.methods) {
            FunctionType declaration = FunctionType.METHOD;
            methodNames.add(method.name.lexeme);
            System.out.println("Add to class " + stmt.name.lexeme + " property: " + method.name.lexeme);

            if (method.name.lexeme.equals("init")) {
                declaration = FunctionType.INITIALIZER;
            }
            resolveFunction(method, declaration);
        }
        \end{verbatim}

        Then, we go into the initializer of the class to track the available fields:
        \begin{verbatim}
    private void resolveFunction(Stmt.Function function, FunctionType type) {
        FunctionType enclosingFunction = currentFunction;
        currentFunction = type;

        if (currentFunction == FunctionType.INITIALIZER && currentClassName != null) {
            System.out.println("looking for this. = <whatever> in constructor.");
            for (Stmt stmt : function.body) {
                // cursed string parsing
                if (stmt instanceof Stmt.Expression) {
                    if (((Stmt.Expression) stmt).expression instanceof Expr.Set) {
                        Expr object = ((Expr.Set) ((Stmt.Expression) stmt).expression).object;
                        Token name = ((Expr.Set) ((Stmt.Expression) stmt).expression).name;
                        if (object instanceof  Expr.This) {
                            classProperties.get(currentClassName).add(name.lexeme);
                        }
//                        System.out.println("object: " + object + " name: " + name);

                    }
                }
            }
        }

        beginScope();
        for (Token param : function.params) {
            declare(param);
            define(param);
        }
        resolve(function.body);
        endScope();

        currentFunction = enclosingFunction;
    }

        \end{verbatim}
        
        The only way
        I was able to implement the basic type checking feature was because I 
        explicilty checked for constructor calls. 
        \begin{verbatim}
    @Override
    public Void visitVarStmt(Stmt.Var stmt) {
        declare(stmt.name);
        if (stmt.initializer != null) {
            resolve(stmt.initializer);
        }

        System.out.println("Declare var");
        if (stmt.initializer instanceof Expr.Call) {
            System.out.println("Initializer is of type Expr.Call");
            Expr callee = ((Expr.Call) stmt.initializer).callee;
            System.out.println("Initializer callee" + callee);
            if (callee instanceof  Expr.Variable) {
                System.out.println("Initializer callee is of type Variable");
                String lexeme = ((Expr.Variable) callee).name.lexeme;
                if (classProperties.containsKey(lexeme)) {
                    System.out.println("Initializer callee is a constructor");
                    // we have found a constructor
                    varClass.put(stmt.name.lexeme, lexeme);
                    System.out.println(stmt.name.lexeme + " is known to be of class " + varClass.get(stmt.name.lexeme));
                }
            }
        }
        define(stmt.name);
        return null;
    }
        \end{verbatim}

        PS. This was genuinely one of the most difficult problems to implement in 
        Java. The lack of pattern matching made it very verbose to check for 
        specific expressions whose type we know at compile time, hence the massive 
        number of `instanceof' calls and casting.
    }
\end{enumerate}

\end{document}
